\scanpage{205}
Ainsi, $\unsure{\Pi|\pi}_1(X)$ n’opère évidemt pas sur $\deleted{\unsure{\bar{\xi}}}\unsure{\xi|\zeta}_{x,\mathbb{Z}}$, ni \unsure{$=$|même} sur $\xi_{x,\mathbb{R}}$, mais l’élément $f_\infty$ de $\Pi_1(X)$ (frobénius infini) opère sur $\xi_{x,\mathbb{Z}}$ et à fortiori sur $\xi_{x,\mathbb{R}}$ ; on le prolonge \uline{linéairement} en $\tau=f_\infty$ \ldots.

\uline{N.B.} La \uline{conjecture de Hodge} interprète la filtration \uline{induite} sur la coh. ordinaire transcendante $\xi_{x,\mathbb{Z}}\otimes_{\mathbb{Z}}\mathbb{Q}\inserted{=\xi_{x,\mathbb{Q}}}$ par la filtration arithmétique ($=$ motivique, par le niveau), en termes de la position \unsure{par.} à la graduation [i.e. en termes de la filtration de Hodge --- non motivique --- et du $f_\infty$]. Elle s’exprime en disant que

\deleted{$\unsure{F|\xi}$ de niveau $\leq\unsure{p|\rho}\iff$}

\[
\begin{aligned}
&\text{La \uline{partie de niveau}}\ \leq p\quad\text{\uline{dans}}\quad\xi_{x,\mathbb{Q}}(F)\\
&\text{\uline{n’est autre que}}\quad \xi_{x,\mathbb{Q}}(F)\cap\xi_{x,\infty}(F)^{\unsure{(p)|p|(\rho)}}\cap f_\infty\xi_{x,\infty}(F)^{(p)}.
\end{aligned}
\]

En particulier, ce dernier sous-espace de $\xi_{x,\mathbb{C}}(F)$ est \uline{motivique} (provient d’un \deleted{\unsure{un}} sous-motif $F^{(p)}$ de $F$), alors que $\xi_{x,\infty}(F)^{(p)}$ et $f_\infty\xi_{x,\infty}(F)^{(p)}$ n’est pas motivique \ldots. \deleted{Si $p=0$, cela donne l’intersection des $2$ réseaux ``\unsure{entiers|rationnels|niveau}''}

% archive sheet number: 204
